\section{Introduction}

Advanced neural recording technologies now enable neuroscientists to capture unprecedented volumes of high-dimensional neural\footnote{\fontsize{8}{8}\selectfont By default, ``neural'' and ``neuron'' refer to neurobiology; e.g. we call brain spike data ``neural activity''. In contrast, activity and neurons from artificial neural networks are explicitly prefaced; e.g. ``model neurons'', ``hidden layer neuron''. For neural data, ``unit'' denotes a signal source that can but does not necessarily imply correspondence to a single neuron.} data at single-cell resolution~\citep{steinmetz_2021_neuropixels2, raducanu_2017_neuroseeker, cai_2016_miniscope, villette_2019_voltage_2p, ouzounov_2017_three_photon, ahrens_2013_lightsheet}. Understanding the computational principles driving neural activity requires extracting interpretable features from this data. Traditional dimensionality reduction techniques, while useful for visualization and basic analyses, often make invalid assumptions about or altogether fail to disentangle distributed representations, and as a result typically cannot provide a mechanistic understanding of the underlying computations~\citep{cunningham_2014_neural_dr, humphries_2021_dr_principles}. Recent work has produced promising latent variable model (LVM) approaches capable of identifying low-dimensional subspaces that can accurately decode aspects of behavior and environment~\citep{song_2025_langevinflow, schneider_2023_cebra, le_2022_stndt, keshtkaran_2022_autolfads, yu_2009_gpfa, macke_2011_plds, gao_2016_pflds, low_2018_mind, jensen_2020_mgplvm, hernandez_2020_vind, kim_2021_plnde, hurwitz_2021_tndm, schimel_2022_ilqrvae, kudryashova_2023_ctrltndm, ye_2023_ndt2, gondur_2024_mmgpvae, pellegrino_2024_slicetca, sani_2024_dpad, pals_2024_smclr_rnn, zhang_2024_mtm, george_2025_simpl, perkins_2025_mint, schmutz_2025_nce}. However, these methods typically value decoding accuracy over latent interpretability, and consequently all have at least one of the following limitations: poorly interpretable latents, lossy mapping to and/or poor reconstruction of neural data from latent space, priors on the latent and/or neural data space, requirements of additional behavioral, environmental, and/or trial-structured neural data, supralinear scaling with respect to dataset size, or high implementational complexity.

We address each of these limitations in NLDisco, which provides a simple, user-friendly library for interpretable latent discovery from high-dimensional neural data. We consider a latent's interpretability in two key aspects: 1) its correspondence to a specific external variable -- a ``natural'' behavioral or environmental feature\footnote{\fontsize{8}{8}\selectfont  A ``feature'' is a sufficiently interpretable latent.}; 2) its explicit composition from contributing neural activity. Unlike other approaches that require a search for meaningful directions or dynamics in a latent space, NLDisco outputs individual latents that can be directly assessed for interpretability. Additionally, while we showcase NLDisco on spike data, it can be readily used with virtually any other neural recording modality as the only data requirement is any predefined spatiotemporal binning.

NLDisco trains hidden layer neurons in shallow, overcomplete, sparse~\citep{olshausen_1997_sparse_overcomplete} encoder-decoder (SED) models~\citep{vincent_2008_denoising_autoencoders} to learn interpretable representations of neural activity. An overcomplete dictionary provides capacity for many candidate features, while sparsity limits how many contribute to any one reconstruction. Together, these encourage explanations of population activity in terms of a set of reusable patterns, making individual latents and their contributions easier to interpret. The overarching field of sparse dictionary learning has had recent successes in mechanistic interpretability, in which models have typically been implemented as sparse autoencoders (SAEs), transcoders (CLTs), or crosscoders (SCCs)~\citep{lindsey_2024_crosscoders, cunningham_2023_saes, bricken_2023_towards_monosemanticity, templeton_2024_scaling_monosemanticity, dunefsky_2024_transcoders, ameisen_2025_circuit_tracing, lindsey_2025_biology_llm}. For neural data, SEDs are additionally motivated by the principle that we need not explicitly impose structure on a latent space nor model its dynamics to discover interpretable features.

Specifically, NLDisco bypasses the need to enforce non-Markovian dynamics, which are often a property of an incomplete observation of a system rather than the system itself; any dynamical system can be represented as Markovian via sufficient state augmentation ~\citep{takens_1981_embedding}. Given this, even methods that appear non-Markovian can be seen as attempts to approximate a complete underlying state. For example, the forward dynamics of a recurrent neural network (RNN) are Markovian in its hidden state ~\citep{sussillo_2013_rnn_dynamics, goodfellow_2016_rnn}, and generalized linear models (GLMs) with history filters ~\citep{pillow_2008_glms, truccolo_2005_pointprocess} use past activity to augment the present state. This perspective aligns with several findings in neuroscience, from activity in motor cortex reflecting a low-dimensional, implicitly Markovian dynamical system ~\citep{churchland_2012_population_dynamics, shenoy_2013_dynamical_perspective}, to predictive-coding formulations that posit that the brain maintains an internal state sufficient to predict sensory streams, i.e., a Markovian latent process ~\citep{rao_1999_predictive_coding, doya_2007_bayesian_brain, friston_2010_free_energy}. Grounded in this perspective, NLDisco forgoes the challenge of modeling intricate dynamics and instead directly extracts meaningful constituent features of the underlying neural state (\autoref{figure:interpretable_latents_vs_latent_space}).
